\documentclass[11pt,a4paper]{article}

\usepackage[margin=1in]{geometry}
\usepackage{amsmath,amssymb,amsfonts}
\usepackage{booktabs}
\usepackage{hyperref}
\usepackage{enumitem}

\hypersetup{
    colorlinks=true,
    linkcolor=blue,
    citecolor=blue,
    urlcolor=blue
}

\title{Synthetic OD Matrix Generation:\\Problem Statement and Constraints}
\author{OD Estimation Project}
\date{\today}

\begin{document}
\maketitle

\begin{abstract}
This note states the problem of generating realistic Origin--Destination (OD) matrices from first principles for machine-learning training. It specifies inputs and outputs, hard and soft transportation constraints, desired statistical properties, and the requirement that samples lie on a low-dimensional latent manifold rather than filling the full $N^2$-dimensional space uniformly.
\end{abstract}

\section{Problem Statement}

\subsection{Objective}

Develop a statistical generator capable of producing an arbitrary number of realistic Origin--Destination (OD) matrices for a transportation network.

Each generated OD matrix should satisfy fundamental transportation principles and network constraints while exhibiting sufficient variability to serve as high-quality training data for machine learning models.

Unlike traditional approaches that perturb an existing OD matrix, the proposed generator should synthesize OD demand \textbf{from first principles}, producing diverse yet realistic travel demand patterns.

\subsection{Inputs}

The generator requires only:

\begin{itemize}
    \item Transportation Analysis Zones (TAZs)
    \item Pairwise distance (or travel cost/time) matrix
    \item Desired total number of trips
    \item Number of OD matrices to generate
\end{itemize}

No historical OD matrix is required.

\subsection{Output}

Generate
\[
\{T^{(1)},T^{(2)},\ldots,T^{(K)}\}
\]
where
\[
T^{(k)}\in\mathbb{R}^{N\times N}
\]
is a valid OD matrix.

\section{Constraints}

\subsection{Non-negativity}

Every OD flow represents trips. Therefore,
\[
T_{ij}\ge 0.
\]

\subsection{No Intrazonal Trips (Optional)}

If intrazonal trips are excluded,
\[
T_{ii}=0.
\]
Otherwise, they may be permitted.

\subsection{Fixed Total Demand}

Every generated matrix should represent the same total daily travel demand:
\[
\sum_{i,j}T_{ij}=G,
\]
where $G$ is a user-specified constant.

\subsection{Production--Attraction Consistency}

The total number of productions must equal the total number of attractions:
\[
\sum_i P_i = \sum_j A_j = G,
\]
where
\[
P_i=\sum_j T_{ij}, \qquad A_j=\sum_i T_{ij}.
\]
If production and attraction vectors are prescribed, every generated matrix must satisfy them exactly. Otherwise, these vectors are generated as part of the synthesis process.

\subsection{Distance Decay}

Trips generally decrease with increasing spatial separation between origin and destination. For two OD pairs with
\[
d_{ij}>d_{mn},
\]
the expected flow should satisfy
\[
\mathbb{E}[T_{ij}]<\mathbb{E}[T_{mn}].
\]
However, this is \textbf{not} a hard constraint. Long-distance trips must still occur with non-zero probability. Thus,
\[
T_{ij}\propto f(d_{ij}),
\]
where $f(d)$ is a monotonically decreasing function such as
\begin{itemize}
    \item exponential decay,
    \item power-law decay,
    \item logistic decay,
    \item gravity-based decay.
\end{itemize}
The distance decay should be interpreted statistically rather than deterministically. While nearby zones generally exchange more trips, some long-distance OD pairs should still exhibit high demand due to latent socioeconomic or transportation factors.

\subsection{Random Corridor Dominance}

Each generated OD matrix should exhibit a unique set of dominant travel corridors. Specifically:
\begin{itemize}
    \item A small fraction of OD pairs should carry very high demand.
    \item Most OD pairs should carry moderate demand.
    \item Many OD pairs should carry low demand.
\end{itemize}
The identity of these dominant corridors should vary from one generated matrix to another. The generator must therefore avoid repeatedly producing the same dominant OD pairs.

\subsection{Heavy-Tailed Flow Distribution}

Cell values should approximately follow a heavy-tailed distribution. Consequently:
\begin{itemize}
    \item Most OD pairs should carry small demand.
    \item A few OD pairs should carry very large demand.
\end{itemize}
This reproduces empirical transportation demand patterns observed in real road networks.

\subsection{Approximate Reciprocity (Symmetry)}

Daily travel demand often exhibits reciprocal movement. Therefore,
\[
T_{ij}\approx T_{ji}.
\]
Exact equality is unnecessary. The generator should preserve only statistical symmetry while allowing directional imbalances.

\subsection{Spatial Coherence}

Nearby OD pairs should exhibit correlated behavior. If one OD pair forms a major travel corridor, neighboring origins or destinations should have an increased probability of generating similar travel demand. Demand should therefore appear spatially coherent rather than as isolated random spikes.

\subsection{Smooth Origin and Destination Importance}

Each origin possesses a production strength. Each destination possesses an attraction strength. Zones with larger production strengths tend to generate more trips; zones with larger attraction strengths tend to receive more trips. These strengths should vary across generated matrices.

\subsection{Diversity}

Generated matrices should not be nearly identical. For two generated matrices $T^{(a)}$ and $T^{(b)}$, require
\[
\|T^{(a)}-T^{(b)}\|_F>\tau,
\]
where $\tau$ is a predefined diversity threshold.

\subsection{Entropy}

The generator should maximize dataset diversity while maintaining transportation realism. Different OD cells should experience varying demand distributions across the generated dataset.

\subsection{Sparsity}

Depending on the intended application, many OD pairs should contain zero, negligible, or low flow. Completely dense OD matrices are generally unrealistic.

\subsection{Network Connectivity}

Trips may only be assigned between connected origin--destination pairs. The generator should never allocate demand to disconnected zones.

\subsection{Statistical Stationarity}

All generated matrices should originate from the same underlying stochastic process. No generated matrix should be an obvious statistical outlier unless explicitly requested.

\subsection{Scalability}

The generator should efficiently produce $10^6$ or more synthetic OD matrices without significant computational overhead.

\section{Desired Statistical Properties}

Across the generated dataset, the following properties should hold:

\begin{itemize}
    \item Constant total demand.
    \item Realistic production distributions.
    \item Realistic attraction distributions.
    \item Heavy-tailed OD flow distribution.
    \item Distance-decay behavior.
    \item Approximate reciprocity.
    \item High entropy.
    \item High pairwise diversity.
    \item Spatial coherence.
    \item Stable marginal statistics.
    \item No duplicate matrices.
    \item Realistic trip-length distribution.
    \item Controllable randomness through a small number of interpretable parameters.
\end{itemize}

\section{Additional Constraint: Low-Dimensional Latent Structure}

Real-world OD matrices are not arbitrary. Instead, they are governed by a relatively small number of latent factors such as
\begin{itemize}
    \item residential population,
    \item employment,
    \item commercial centers,
    \item educational institutions,
    \item industrial zones,
    \item recreational areas,
    \item accessibility,
    \item socioeconomic characteristics.
\end{itemize}
Consequently, although an $N\times N$ OD matrix contains $N^2$ entries, its effective dimensionality is considerably smaller.

A realistic generator should therefore produce OD matrices that lie on a low-dimensional manifold rather than uniformly occupying the entire $N^2$-dimensional space. This property explains why dimensionality reduction techniques such as Principal Component Analysis (PCA) and Autoencoders can effectively compress OD matrices while preserving most of their information.

\section{Overall Problem Definition}

Design a constrained probabilistic generator capable of sampling realistic Origin--Destination matrices from a transportation-aware probability distribution defined by
\begin{itemize}
    \item network topology,
    \item distance-decay effects,
    \item production--attraction consistency,
    \item spatial coherence,
    \item approximate reciprocity,
    \item heavy-tailed corridor demand,
    \item dataset diversity,
    \item scalability.
\end{itemize}
The generated matrices should satisfy transportation constraints while remaining statistically diverse, physically plausible, and suitable for large-scale machine learning applications.

\end{document}
